\section{The coefficient on logarithmic scales}\label{sec:profile}

Proposition~\ref{prop:passage} locates a change in a torus variance at
logarithmic scale $T$. We now identify the free and pinned height laws
on either side of that scale. The comparison estimates are finite-graph
statements, whereas the height limit theorems in \cite{X} are stated at
fixed inverse temperature. We therefore formulate the limit argument
for a sequence of temperatures and retain the order of its auxiliary
limits.

Throughout this section let $b_k<\bc$, $b_k\to\bc$, and put
$\delta_k=\sqrt{\bc-b_k}$. Passing to subsequences will always preserve
these notations. All finite comparison constants may be chosen for one
fixed compact interval of positive temperatures containing the $b_k$
and $\bc$.

\subsection{Extraction of the coefficient profile}

Let $B_m=\{0,\ldots,m-1\}^2$, let $A_{B_m}$ be its free unit-edge
Laplacian, and define
\begin{equation}\label{eq:profile-coefficient}
 g_m(x)=x_1/m,\qquad
 a_m(b)=\Var_{B_m,b}^{\mathrm{free}}(h,A_{B_m}g_m).
\end{equation}
The test in \eqref{eq:profile-coefficient} is the difference of the
right and left face averages. The two coordinate affine tests have
covariance matrix $a_m(b)I_2$, by the reflections and quarter turns of
the square. Their reference costs are bounded, so $a_m(b)\le C$
uniformly in $m$ and in the temperatures under consideration.

Cutting an $M$-square into $m$-squares and strips of width less than $m$
gives the comparison
\begin{equation}\label{eq:profile-cut}
 \sqrt{a_M(b)}\le \sqrt{a_m(b)}
       +C\bigl(m^{-1}+m/M\bigr)^{1/2},\qquad M\ge m.
\end{equation}
Indeed, omit the seam and strip flows in the large square's affine
test. Their squared reference cost is $O(m^{-1}+m/M)$. Releasing the
cell constraints increases covariance and makes the retained cell
tests independent. The triangle inequality in $L^2$ and the reference
moment bound then give \eqref{eq:profile-cut}. This is \cite[Equation~(2.6)]{X}; its constant is uniform here because it uses only the
compact-temperature reference bound.

\begin{proposition}[Extraction on logarithmic scales]\label{prop:profile}
Every sequence $b_k\to\bc$ from below has a subsequence and a
nonincreasing function $\mathcal A:(0,\infty)\to[0,a]$, where
$a=8\pi$, with the following property. At each continuity point
$w>0$ of $\mathcal A$, every integer sequence $m_k$ satisfying
$\delta_k\log_L m_k\to w$ obeys
\begin{equation}\label{eq:profile-limit}
                         a_{m_k}(b_k)\longrightarrow\mathcal A(w).
\end{equation}
One may choose $\mathcal A$ to be right continuous.
\end{proposition}

\begin{proof}
For positive rational $q$, set $m_k(q)=\lfloor L^{q/\delta_k}\rfloor$.
A diagonal subsequence gives a limit of $a_{m_k(q)}(b_k)$ for every
such $q$. If $q<q'$, then $m_k(q)/m_k(q')\to0$ and
$m_k(q)\to\infty$. Equation~\eqref{eq:profile-cut} orders the two
limits. Denote these rational limits by $\alpha(q)$ and define their
right-limit regularization by
\[
             \mathcal A(w)=\sup_{q>w,\ q\in\mathbb Q}\alpha(q).
\]
This is nonincreasing and right continuous. Values at rational jump
points need not be preserved; only continuity points will be used.

For $q_-<w<q_+$, a sequence as in the statement eventually lies between
$m_k(q_-)$ and $m_k(q_+)$. Applying \eqref{eq:profile-cut} in these two
directions bounds the limsup and liminf of $\sqrt{a_{m_k}(b_k)}$ by
the values at $q_-$ and $q_+$. Letting these rationals tend to the
continuity point $w$ proves \eqref{eq:profile-limit}.

Centered covariance increases with $b$, as follows either from the
chain approximation or from the finite comparisons of
Section~\ref{sec:finite}. Thus $a_m(b_k)\le a_m(\bc)$. The critical
height coefficient is $a$ by \cite[Section~6]{X}, and
$a_m(\bc)\to a$ by \cite[Theorem~2.4]{X}.
Consequently $0\le\mathcal A\le a$.
\end{proof}

Fix for the moment a continuity point $w$, a sequence $n_k$ with
$\delta_k\log_Ln_k\to w$, and write $a_*=\mathcal A(w)$.
Multiplying $n_k$ by any fixed positive constant, or rounding it by a
bounded number of lattice steps, leaves this logarithmic scale
unchanged. This observation supplies the coefficient limit on all
fixed macroscopic pieces used below.

\subsection{Free Gaussian limits at a moving temperature}

We use scaled densities with the following convention. A lattice
profile $f_k$ represents $f\in L^2(U)$ if its piecewise constant
scaled density converges to $f$ in $L^2$. For a fixed bounded piecewise
continuous $f$ one may take $f_k(x)=n_k^{-2}f(x/n_k)$, using compatible
cell averages if needed along its discontinuities. When $\int_Uf=0$,
make the total exactly zero by replacing this profile with
$f_k-(\sum_xf_k(x))u_{U_k}$, where $u_{U_k}$ is uniform probability.
The subtracted total tends to zero. Smooth shifts $g$ are sampled as
$g_k(x)=g(x/n_k)$. A boundary probability is normalized by the number
of vertices on that face. Joint statements below include every fixed
finite collection of these observations.

\begin{lemma}[Free limits in a fixed rectangle]\label{lem:profile-free}
Let $U_k$, rescaled by $n_k^{-1}$, be free rectangles converging to a
fixed nondegenerate axis rectangle $U$. For smooth shifts defined near
$\overline U$,
\begin{equation}\label{eq:profile-free-gradient}
 \log\E_{U_k,b_k}^{\mathrm{free}}e^{(h,A_{U_k}g_k)}
       \longrightarrow \frac{a_*}{2}\int_U|\nabla g|^2.
\end{equation}
The joint limits are centered Gaussian and all joint moments converge.
For zero-total density smears the limiting covariance is
$a_*(-\Delta_{U,N})^{-1}$ on mean-free functions. Face probabilities
minus volume probabilities of the same total are also allowed.
These conclusions include $a_*=0$.
\end{lemma}

\begin{proof}
We give the moving-temperature version of the cutting and folding
argument in \cite[Theorem~2.4]{X}. Its one
infinite-size affine comparison is replaced below by a finite
macroscopic enclosure.

First consider an affine shift $g(x)=p\cdot x$. Tile most of $U$ by
squares of a small fixed macroscopic side, leaving strips of
arbitrarily small total area. The coefficient of every retained
square tends to $a_*$ by \eqref{eq:profile-limit}. Releasing the cuts,
using independence of the resulting laws, and bounding the strip
flows by reference cost give
\begin{equation}\label{eq:profile-affine-upper}
 \limsup_k\Var(h,A_{U_k}g_k)\le a_*|U|\,|p|^2.
\end{equation}
The strip bound is sent to zero after $k\to\infty$.

To prove the reverse inequality, place $U$ strictly inside a fixed
square $S$ and partition $S\setminus U$ into finitely many rectangles
of positive width. Use the actual $U_k$ endpoints for the complementary cuts; their
rescaled lengths converge to the prescribed positive lengths.
The full affine observation on $S_k$ is the sum of the affine
observations on these pieces and $U_k$, apart from seam flows of
reference cost $o(1)$. Its variance tends to $a_*|S|\,|p|^2$.
Cutting increases covariance; each piece has the upper bound
\eqref{eq:profile-affine-upper}. The sum of those upper bounds is
exactly $a_*|S|\,|p|^2$. Hence the bound must be attained on each
piece. Otherwise a subsequence with a fixed deficit on one piece,
together with the upper bounds on the others, would contradict the
limiting variance on $S_k$. This proves affine saturation on $U$.
All enclosing lengths are asymptotic to fixed positive multiples of $n_k$.

For a smooth or continuous piecewise smooth $g$, approximate its
gradient on small fixed macroscopic rectangles by constant vectors.
The seams have vanishing reference cost; the interior error is bounded
by its squared $L^2$ gradient error. Letting the rectangles shrink
after the mesh limit gives
\begin{equation}\label{eq:profile-partial-upper}
 \limsup_k\Var(h,A_{U_k}g_k)\le a_*\int_U|\nabla g|^2.
\end{equation}
The same estimate holds for partial gradient observations restricted
to polygonal pieces. Explicitly, for a piece $H$ the partial test is
\[
 J_{H,k}(g)=\sum_{\substack{\{x,y\}\in E(U_k)\\x,y\in H_k}}
              (g_k(y)-g_k(x))(h_y-h_x),
\]
where $H_k=\{x\in U_k:x/n_k\in H\}$, and its energy bound is
$a_*\int_H|\nabla g|^2$. Edges crossing a
fixed cut contribute a vanishing reference-cost error.
Call equality \emph{saturation}. On a convergent
joint covariance subsequence, the difference between the quadratic
form on the right of \eqref{eq:profile-partial-upper} and the limiting
covariance is positive semidefinite. Its null space is linear.
Saturation therefore passes to linear combinations, to
$L^2$-gradient limits, and to restrictions obtained by cutting an
enclosing rectangle and comparing the sum of the energy bounds.

For completeness, the finite reflection step used to generate
nonaffine saturated shifts has no temperature-limit hypothesis.
Work first in a rectangle symmetric about the reflection line.
Choose the representative with one height on the line fixed at zero,
and let $\mathcal T$ be the sigma-field generated by its full trace.
For an odd saturated shift, denote by $Y_+$ and $Y_-$ the partial
tests on the two open sides after transporting exceptional
coefficients to the line in a reflection-symmetric manner.
Suppose that the off-line coefficients of $Y_+$ are nonnegative;
those of $Y_-$ are then nonpositive. The subdivided quadratic law is
conditionally associated, and this property passes to the primary
graph. Consequently
$\Cov(Y_+,Y_-\mid\mathcal T)\le0$.
Reflection fixes the prescribed trace and exchanges the two sides,
while oddness changes the sign of the test. Thus the conditional
means satisfy the exact identity
\[
 \E(Y_-\mid\mathcal T)=-\E(Y_+\mid\mathcal T).
\]
The total covariance is therefore
\[
 \Cov(Y_+,Y_-)
   =\E\Cov(Y_+,Y_-\mid\mathcal T)
       -\Var\bigl(\E(Y_+\mid\mathcal T)\bigr)\le0.
\]
The transport has vanishing reference cost, so its contribution to
covariances tends to zero by the uniform second-moment bound.
The saturated sum and the separate energy bounds
\eqref{eq:profile-partial-upper} then force both original partial
tests to saturate. An odd shift vanishes on its reflection
line, so the weak gradient of $g\one_H$ is
$\one_H\nabla g$. Its lattice gradient test differs from
$J_{H,k}(g)$ only by the stated vanishing-cost seam. Thus the folded
shift itself saturates.

Here are the shifts and transports from
\cite[proof of Theorem~2.4, Equation~(2.8)]{X}. For a hinge at $t$,
use a fixed enclosing rectangle symmetric about $x_1=t$ and the odd
affine seed $g=x_1-t$. Transport the negative coefficients adjacent to
that vertical cut to the cut. There are $O(n_k)$ of them,
each of size $O(n_k^{-1})$, on bounded-length paths of bounded overlap.
Their total squared cost is $O(n_k^{-1})$. This yields the hinges
$(x_1-t)_+$; restriction gives their saturation on the target
rectangle. Linear combinations and energy approximation then give
every smooth function of $x_1$, and likewise of $x_2$. For a diagonal cut write $Y=x_2+t$ and start with
$F(x_1)-F(Y)$ on a square symmetric under interchanging $x_1,Y$.
Initially impose $F'''<0$ and positive derivatives at the two end
faces. The interior coefficients in $x_1>Y$ are nonnegative because
the centered second differences of $F$ decrease with their center;
the two outer-face coefficients have positive leading terms.
Exceptional diagonal-strip coefficients again have total transport
cost $O(n_k^{-1})$. We obtain saturation of
\begin{equation}\label{eq:profile-fold}
             (F(x_1)-F(x_2+t))\one_{\{x_1>x_2+t\}}.
\end{equation}
Every smooth $F$ is a difference of two functions with those strict
properties, by adding a sufficiently large cubic of negative third
derivative and then a sufficiently large positive linear function.
Thus \eqref{eq:profile-fold} holds for arbitrary smooth $F$.
Rounding a reflection line changes the squared gradient cost by
$O(n_k^{-1})$.

Given smooth $g$, choose a bounded interval $I$ containing all
$x_1-x_2$ in the target rectangle. Choose a smooth family $F_t$,
compactly supported in $I$ after slightly enlarging $I$, with
$F_t'(x)=-g_{12}(x,x-t)$ whenever $(x,x-t)$ lies in that rectangle.
Integrate \eqref{eq:profile-fold} over $t\in I$. Indeed, differentiating first in
$x_1$ and then in $x_2$ gives the distributional identity
\[
 \partial_{12}\int_I
 (F_t(x_1)-F_t(x_2+t))\one_{\{x_1>x_2+t\}}\,dt
          =-F'_{x_1-x_2}(x_1)=g_{12}(x_1,x_2).
\]
The difference from $g$ is a sum of one-variable functions. Riemann
sums converge in gradient energy, including across the moving
reflection lines. Linearity and energy approximation therefore prove
saturation for every smooth $g$.

The finite centered comparison gives
$\log\E e^X\ge\Var(X)/2$, providing the Laplace lower bound.
For the upper bound, cut into cells of fixed macroscopic side $r$.
A retained smooth-gradient test $X$ on one cell has reference-cost
square root $O(r)$. The normalized exponential moment estimate gives,
for each fixed $p>1$,
\[
 \log\E e^{pX}=\tfrac12p^2\Var(X)+O_p(r^3),
 \qquad \E(|X|^3e^{p|X|})=O_p(r^3),
\]
with constants uniform in $k$. H\"older's inequality and the same
reference bound remove the seam errors. There are $O(r^{-2})$ cells.
First send $k\to\infty$, then $r\downarrow0$, and finally
$p\downarrow1$. This proves \eqref{eq:profile-free-gradient}.
Apply it to real linear combinations to obtain joint Gaussian limits;
the exponential bounds supply convergence of moments.

For densities, approximate by finite sums of Neumann eigenfunctions
and use their smooth inverse-Laplacian shifts. The discrete Neumann
form approximation and Poincar\'e inequality bound the reference cost
of the error by its squared scaled $L^2$ norm. Using the actual
rectangle endpoints before passing to their limits handles the
rounding. Finally move a face probability to a parallel strip of
width $r$. Parallel normal flows cost $O(r+n_k^{-1})$; the strip is a
density smear. Let $k\to\infty$ and then $r\downarrow0$.
\end{proof}

\subsection{Fractional means}

For a free graph and a profile $s$ of integer total, write
$\phi(s)=\E e^{i(h,s)}$. It is independent of the representative of
$h$ modulo a common $2\pi$ translation. Finite comparison gives
$0<\phi(s)\le1$. If $v$ has zero total, the logarithmic Hessian bound
and the fact that $\log\phi\le0$ imply
\begin{equation}\label{eq:profile-phase-transfer}
 \left|\sqrt{-\log\phi(s+v)}-\sqrt{-\log\phi(s)}\right|
                  \le\sqrt{\Var(h,v)/2}.
\end{equation}
Indeed, for all real $t$ let $F(t)=\log\phi(s+tv)$ and
$V=\Var(h,v)$. The inequalities $F''\ge-V$ and $F\le0$ give
$|F'|^2\le-2VF$ by optimizing the quadratic lower bound for $F$;
integration gives \eqref{eq:profile-phase-transfer}, as in
\cite[Equation~(2.2)]{X}.

\begin{lemma}[Uniform fractional means]\label{lem:profile-phase}
Under the hypotheses of Lemma~\ref{lem:profile-free}, suppose $a_*>0$.
The fractional mean modulo $2\pi$ converges to uniform measure,
independently of the Gaussian observations in that lemma. For every
fixed nonzero integer $l$ and every zero-total profile $v_k$ of bounded
reference cost,
\begin{equation}\label{eq:profile-phase-zero}
 \E\exp\{i(h,lu_{U_k}+v_k)\}\longrightarrow0,
\end{equation}
where $u_{U_k}$ is uniform probability on $U_k$.
\end{lemma}

\begin{proof}
We describe the soft joining and replica argument of
\cite[Theorem~2.5, Equations~(2.10)--(2.13)]{X}, keeping
its lattice and geometric limits separate. Start with a square of side
$m=m_k$ on the chosen logarithmic scale. Let $u_m$ be its volume
probability, $\rho_m$ its right-face probability, and set
\[
 X=(h,A_{B_m}(x_1/m)),\qquad Y=(h,\rho_m-u_m).
\]
Lemma~\ref{lem:profile-free} gives joint Gaussian limits, with
$\Var X\to a_*$ and $\Cov(X,Y)\to a_*/2$.
Suppose along a subsequence that $\phi(lu_m)$ is bounded away from
zero for a nonzero integer $l$. Uniform exponential bounds provide a
further subsequence on which
\[
 F_m(t,v)=\E e^{i(h,lu_m)+itX+ivY}
\]
converges locally with all derivatives. On real arguments the limit
$F$ is strictly positive by \eqref{eq:profile-phase-transfer}.

The replica identity is a finite quadratic-lattice identity. Let
$F_{m,N}$ denote the preceding transform with each primary edge
replaced by $N$ Gaussian chain bonds and the observations evaluated
at the primary endpoints. For two independent fields under this
subdivided law put $S=h_1+h_2$, $D=h_1-h_2$.
Choose both representatives to vanish at the same reference vertex.
At every vertex of the subdivided graph, $S\equiv D$ modulo
$4\pi\Z$. Conditional on this full common parity class, $S$ and $D$
are independent and identically distributed. If $V_{ij}$ is the conditional second moment
of $S_i,S_j$, for linear observations indexed by $i,j$, then
\begin{equation}\label{eq:profile-replica}
 \E V_{ij}=2\Cov(h_i,h_j),\qquad
 \Cov(V_{ij},V_{rs})=2\kappa(h_i,h_j,h_r,h_s),
\end{equation}
where all moments in this identity are in the subdivided law and
$\kappa$ is its fourth joint cumulant. This follows by expanding
$\E S_iS_jD_rD_s$ in the original two copies. Differentiating
$F_{m,N}$ in its observation variables gives
\[
 \E S_iS_j e^{i(l(h_1-h_2,u_m)+tX(D)+vY(D))}
       =-2(F_{m,N}F_{m,N,ij}-F_{m,N,i}F_{m,N,j}).
\]
In the conditional expectation on the left, replace $V_{ij}$ by its
mean. The error is at most $\sqrt{\Var V_{ij}}$. First take the chain
limit on each fixed primary graph. Then the Gaussian moment limits
in Lemma~\ref{lem:profile-free} make the needed fourth cumulants zero.
For $ij=XX,XY$ this yields
\[
 (\log F)_{tt}=-a_*,\qquad (\log F)_{tv}=-a_*/2.
\]
Reflection gives $F_t(0,0)=0$. Consequently
\begin{equation}\label{eq:profile-phase-transform}
 F(t,v)=F(0,v)\exp\{-a_*t^2/2-a_*tv/2\}.
\end{equation}

Softly join two horizontal squares as follows. Keep independent
seam increments with the Bessel law and put
$H_2=(h_1,\rho_R)-(h_2,\rho_L)-\bar\eta$,
where $\bar\eta$ is their mean with the joining orientation.
Penalize this mismatch by $e^{-H_2^2/2}$. Sum relative cell translations; quotient by their
common translation. Adding exact seam matching recovers the joined
rectangle. The soft variance of the sum of the internal affine tests
is bounded above by twice the free square variance and below by its
variance in the joined rectangle. The joined full affine test differs
from this sum by a seam flow of squared cost $O(m^{-1})$.
Lemma~\ref{lem:profile-free} therefore makes the soft variance tend to
$2a_*$.

The comparisons just used with summed translations can be obtained
from ordinary finite comparisons without reweighting the original
free cells. First confine each cell mean by an auxiliary precision
$\zeta\sum_i\bar h_i^2$. For each representative mean $\mu$,
its translation sum is
$\sum_{q\in\Z}e^{-\zeta(\mu+2\pi q)^2/2}$.
After division by its Gaussian integral scalar, this tends uniformly
to one as $\zeta\downarrow0$. After adding the soft penalty, divide
out the common-translation scalar; the relative-translation Gaussian
sum is uniformly bounded in its representative-dependent shift.
The finite precision comparisons at $\zeta>0$ therefore pass to
$\zeta=0$ by dominated convergence, using the uniform internal
moments. This justifies both variance comparisons with the original
free-cell product law.

Poisson summation of the relative translations expresses the soft
characteristic transform as a positive sum of integer sectors, with
weights $e^{-q^2/2}\chi_m(q)$; here $0<\chi_m(q)\le1$ is the
independent seam characteristic. Explicitly,
$\chi_m(q)=\exp\{mb(\cos(2\pi q/m)-1)\}$, so
$\chi_m(q)\to1$ for fixed $q$, uniformly in the temperatures here.
To specify the sector functions, write
$F_m^{(q)}(t,v)=\E e^{iq(h,u_m)+itX+ivY}$. If $\rho_R,\rho_L$
are the two vertical face probabilities, reflection and charge
reversal give
\[
 \psi_q(t)=\phi(q\rho_R+tA(x_1/m))
               \phi(-q\rho_L+tA(x_1/m))
           =F_m^{(q)}(t,q)^2.
\]
Each sector satisfies $(\log\psi_q)''\ge-2\Var X$. The logarithmic Hessian of a positive
sum is the mean of the sector logarithmic Hessians plus the variance
of the sector logarithmic first derivatives. Thus the soft variance
is at most twice the free variance minus this latter, nonnegative
variance. Differentiation under the sector sum is justified uniformly:
with $V=2\Var X$,
\[
 \psi_q| (\log\psi_q)'|^2
       \le2V(-\psi_q\log\psi_q)\le2V/e,
\]
and raw second derivatives are bounded by $V$. The Gaussian weights
are summable.

At $t=0$ the normalizing constant is
$Z_m=\sum_qe^{-q^2/2}\chi_m(q)\psi_q(0)$, and
$1\le Z_m\le\sum_qe^{-q^2/2}$ because the zero sector is one.
Its normalized weight $p_0$ is therefore bounded below and its
logarithmic derivative is zero. The sector $q=l$ also has normalized
weight $p_l$ bounded below: \eqref{eq:profile-phase-transfer} moves $lu_m$ to a
seam face at bounded cost. By reflection, charge reversal and
\eqref{eq:profile-phase-transform}, its logarithmic derivative tends
to $-a_*l$, up to orientation. If $D_q$ denotes the sector
logarithmic derivative, their contribution bounds the variance of
derivatives below by $p_0p_l(D_l-D_0)^2$. Its liminf is strictly
positive, contradicting convergence of the
soft variance to $2a_*$. Hence $\phi(lu_m)\to0$.
Equation~\eqref{eq:profile-phase-transfer} then proves
\eqref{eq:profile-phase-zero} for square profiles of bounded cost.

For a fixed rectangle, enclose it in a fixed-multiple square. Regard
its mean phase as a phase on the square. The difference from the
square mean has bounded reference cost, so its square coefficient
tends to zero. Cutting the complement decreases characteristic
coefficients and leaves the rectangle phase on one free component;
hence the rectangle coefficient also tends to zero. Bounded-cost
phase transfers inside the rectangle give the stated generality.
Mixed real characteristic functions with the mean-free observations
now vanish whenever $l\ne0$, and for $l=0$ they have the Gaussian
limit. Fourier approximation on the circle proves joint independence.
Uniform exponential moments permit passage with any fixed real tilt
and polynomial insertion in the mean-free observations.
\end{proof}

\subsection{Primary pins at a moving parameter}

We next transfer the free limits to a hard pin.  This requires more
than convergence of finitely many unconditioned observations: a pin
specifies every primary height in a region.  We use the band and cable
argument of \cite[Section~3]{X}, stating its finite estimates explicitly
and verifying the limits with the parameter varying along the sequence.
At a positive continuity value of the coefficient profile, the free
and phase assumptions below are supplied by
Lemmas~\ref{lem:profile-free} and~\ref{lem:profile-phase}.

\begin{lemma}[Moving-parameter pinned limits]\label{lem:profile-pins}
Let \(b_k\) lie in a compact interval \(J\subset(0,\infty)\), and let
\(n_k\to\infty\).  Suppose that, on every fixed nondegenerate rectangle
with mesh \(n_k^{-1}\), the free Bessel height laws at \(b_k\) have
Gaussian Laplace and moment limits for mean-free density smears and
energy-approximable weak Laplacian tests, with coefficient \(a_*>0\).
Suppose also that the fractional mean modulo \(2\pi\) converges jointly
with those tests to an independent uniform variable.

Let \(U\) be a fixed rectangle and let \(p\subset U\) be a nonempty
finite union of aligned positive-area boxes.  Assume that the resulting
pieces have nondegenerate faces and positive-width passages; bounded
lattice rounding is allowed.  Pin every primary vertex of \(p\) to zero.
Define
\[
 \mathcal H_p=\{v\in H^1(U):v=0\text{ almost everywhere on }p\},
 \qquad \mathcal E_U(v,w)=\int_U\nabla v\cdot\nabla w,
\]
and let \(G_p\) be the inverse of this form: for a bounded density \(f\),
\(G_pf\in\mathcal H_p\) satisfies
\(\mathcal E_U(G_pf,v)=\int_Ufv\) for \(v\in\mathcal H_p\).
For every fixed finite collection of bounded piecewise continuous
densities \(f_1,\ldots,f_d\), their primary smears converge jointly,
with all Laplace transforms and polynomial moments, to the centered
Gaussian vector with covariance
\begin{equation}\label{eq:profile-pinned-covariance}
                  a_*\int_U f_iG_pf_j.
\end{equation}
The densities need not have zero integral.
\end{lemma}

\begin{proof}
Write \(n=n_k\) and \(b=b_k\) in this proof, and index primary vertices
by \(x\in nU\cap\Z^2\), with the permitted boundary rounding.
For a continuum density \(f\), put \(f_n(x)=n^{-2}f(x/n)\);
equivalent cell quadratures give the same limits.  All discrete
pairings use ordinary sums: \((h,f_n)=\sum_xh_xf_n(x)\), whereas
\(g_n(x)=g(x/n)\).  Probability averages, including cell averages,
use weights equal to the reciprocal of their full primary counts.
The unit-edge form is denoted by \(A_{U,n}\); its energy tends to
\(\mathcal E_U\).  All unspecified finite-graph constants below are
uniform for \(b\in J\).  We first prove the assertion for
\[
                     X_{g,n}=(h,A_{U,n}g_n),
\]
where \(g\) is smooth on the rectangle, may meet its free walls, and
vanishes in a neighborhood of \(p\).  The final approximation will use
the form grounded on \(p\) to reach arbitrary densities.

Fix a closed neighborhood \(V\) of \(\supp g\), separated from \(p\),
so that the primary support of \(A_{U,n}g_n\) lies in \(V\) for all
sufficiently large \(k\).  Choose a positive-width band \(K\)
separating \(p\) from \(V\), with positive gaps on both sides.
Such a band can be chosen as a finite union of rectangular strips, clipped at the free
walls.  Partition it into rectangular cells \(d\) of side comparable
to \(\ell\), with uniformly bounded aspect ratios.  If \(h_d\) is the
probability average in \(d\), set
\begin{equation}\label{eq:profile-band}
 W_{\ell,\lambda}(h)
   =\exp\left\{-\frac\lambda2\sum_d|d|h_d^2\right\},
 \qquad
 \widehat W_{n,\ell,\lambda}
   =\frac{W_{\ell,\lambda}}{\mathcal P_n(W_{\ell,\lambda})}.
\end{equation}
Here \(\mathcal P_n\) is the sigma-finite free measure obtained by
normalizing the law modulo common \(2\pi\Z\)-translations and then
summing all translations.  The band confines its constant, so the
denominator in \eqref{eq:profile-band} is finite and positive.  We write
\(\E_W\) for integration against \(\widehat W\,d\mathcal P_n\).
Throughout the proof the order is
\begin{equation}\label{eq:profile-band-order}
                k\to\infty,\qquad \ell\downarrow0,
                \qquad \lambda\to\infty.
\end{equation}
In particular, every band partition and strength is fixed before the
moving-parameter mesh limit.

We will prove that these successive band limits give variance
\(a_*\mathcal E_U(g,g)\) for \(X_{g,n}\), and that this value is
no larger than \(\liminf_k\Var_pX_{g,n}\). The free variance upper
bound will then give equality. The comparison is the main task.
The finite exploration argument below controls the linear tilts
introduced by the band and proves that exploration from \(p\) reaches
\(V\) with vanishing probability. When it does not reach \(V\),
the exposed test is zero and the centered posterior variance is at
most the hard-pin variance; the bounded band fourth moment controls
the contribution when it does reach \(V\). We finish by squeezing
Laplace transforms and using the grounded form to approximate density
smears.

\emph{Band limits.}
At a fixed partition and strength, the assumed joint free limits turn
translation summation into integration of the constant against
\(dc/(2\pi)\).  This passage also holds with every fixed polynomial
insertion.  Indeed, if \(S\) is the probability average over \(U\) and
\(\bar h_K=|K|^{-1}\sum_d|d|h_d\), then
\[
 W_{\ell,\lambda}\le
       e^{-\lambda|K|\bar h_K^2/2},
 \qquad S-\bar h_K\text{ has bounded primary reference cost}.
\]
The reference-cost moment and fixed-linear-exponential bounds, uniform
on \(J\), control this difference; the displayed Gaussian factor
controls the sum over translations.  They give uniform integrability
of the normalizing constants and polynomial insertions.  The limiting
law is therefore the Gaussian with precision
\(a_*^{-1}\mathcal E_U+\lambda\sum_d|d|v_d^2\).

The subsequent cell refinement is a statement about these Gaussian
forms.  Integrating their constant first leaves the covariance
\(C_\ell\) of the cell projection of the mean-free Neumann field,
after subtracting its \(K\)-average.  The inverse Neumann Laplacian
on a rectangle is Hilbert--Schmidt, so
\(\sup_\ell\|C_\ell\|_{\mathrm{HS}}<\infty\).  For \(q\ge1\),
if \(t_j\) are the eigenvalues of \(\lambda C_\ell\), the determinant
factor in the normalized \(q\)-th weight moment satisfies
\[
 \frac12\sum_j\{q\log(1+t_j)-\log(1+qt_j)\}
                  \le C_q\sum_jt_j^2.
\]
The constant integral is a bounded positive factor at fixed \(\lambda\).
Poincar\'e's inequality, with the constant controlled by the band,
also bounds every moment of \(S\) in the resulting Gaussian law.
Consequently, for every fixed \(q\ge1\), \(r\ge0\), and \(\lambda>0\),
\begin{equation}\label{eq:profile-band-density}
 \limsup_{\ell\downarrow0}\limsup_{k\to\infty}
 \mathcal P_n\bigl(\widehat W^q(1+|S|)^r\bigr)<\infty.
\end{equation}
Strong convergence of cell projections and the ordinary form
approximation give the band covariance
\begin{equation}\label{eq:profile-band-Gaussian}
 T_\lambda^{-1},\qquad
 T_\lambda=-\Delta_U/a_*+\lambda\one_K.
\end{equation}
These are the ingredients of the band-density argument in
\cite[Lemma~3.3]{X}, now deduced from the assumed sequence limits.
Since \(g\) vanishes near \(K\),
\(T_\lambda(a_*g)=-\Delta_Ug\) in the weak Neumann sense.  Thus the
limiting variance of \(X_{g,n}\) is \(a_*\mathcal E_U(g,g)\).
Its limiting fourth moment is bounded uniformly for
\(\lambda\ge\lambda_0>0\), by Gaussian covariance monotonicity.

\emph{Finite cable estimates.}
Metric-graph Gaussian sign clusters are developed in
\cite[Section~3, Lemma~3.2]{Lupu}; the Brownian-bridge extension of
integer-restricted Gaussian fields is described in
\cite[Section~7]{AHPS}. We use this construction on the finite
Gaussian-chain approximants.
We record the finite inputs used to remove the band and impose the
primary pin.  Replace each primary edge by \(N\) lattice-Gaussian
bonds with increment weight \(q_N^{j^2}\), where \(q_N=b/(2N)\),
and attach the corresponding independent Brownian bridges conditional
on their endpoint heights.  Let \(\mathcal P_{n,N}\) be the resulting
unweighted sigma-finite cable measure, and let \(\E_{W,N}\) and
\(\Pbb_{W,N}\) denote its normalized band law; write
\(\widehat W_N=W/\mathcal P_{n,N}(W)\).  A cable bond is open
when its bridge does not hit zero.  Under the unweighted cable measure,
conditional on the magnitudes and the open graph, the nonzero clusters
have independent fair signs.  Exploration stops
a bridge at its first zero or at its next vertex.  The centered
unvisited law has additional zero pins and increased precision on
the stopped segments.  These finite disintegration and comparison
facts, together with the estimates below, are those of
\cite[Lemmas~2.2, 3.4, and~3.5]{X}.
All cluster quantities in the following argument refer to these
finite-\(N\) measures.  Subdivision is removed by a
\(\limsup_{N\to\infty}\) at each fixed primary graph, before taking
\(k\to\infty\).  Primary observables and their band normalizers
have actual fixed-graph limits, equal to those of the Bessel law.
No limit of the cable clusters is required.  A remainder denoted by
\(o_k(1)\) in a cable estimate means that its absolute value has
\(\limsup_{k\to\infty}\limsup_{N\to\infty}\) equal to zero at the
fixed auxiliary parameters then under consideration.

The dependence on \(b\) is controlled by the whole-chain estimate
\begin{equation}\label{eq:profile-chain-mgf}
 \log\frac{\sum_{j\in\Z}q_N^{j^2}e^{2\pi tj}}
               {\sum_{j\in\Z}q_N^{j^2}}
       \le \frac{C_{J,R}t^2}{N},\qquad |t|\le R,
\end{equation}
for all sufficiently large \(N\), uniformly on \(J\).
This is \cite[Equation~(3.5)]{X} with its parameter dependence made explicit.
Grouping opposite integers proves this directly: the numerator
increment is bounded by a constant times
\(t^2\sum_{j\ge1}q_N^{j^2}j^2e^{2\pi Rj}=O_{J,R}(t^2/N)\).
A primary flow extended through all \(N\) bonds therefore has a
uniform moment cost.  Additional zero pins and increased precision
decrease this centered cost.  No bare resistance of the subdivided
network is used.

For clarity, write \(\mathcal R(f)\) for the minimum squared
unit-edge flow cost of a mean-free primary test \(f\).  After exposing
a cluster of scaled diameter at least \(r>0\), the full-component
unvisited mean has grounded reference cost at most \(V_r\).  If
\(\mathcal D\) is the collection of band cells touched by an exposed
path and \(\ell<r/C\), the remaining averages satisfy
\begin{equation}\label{eq:profile-path-cost}
 \mathcal R_{\rm post}
       \left(\sum_{d\in\mathcal D}|d|c_dh_d^U\right)
          \le C\ell^2\sum_{d\in\mathcal D}|d|c_d^2.
\end{equation}
Here the superscript \(U\) denotes the unvisited contribution,
extended by zero on visited primary vertices.  The same bound holds
for exploration from a filled region.  In particular, the remaining
probability average in a cell reached by a path from a fixed multiple
of its diameter away has uniformly bounded centered posterior moments.
These statements follow from path Poincar\'e on the unit-edge network:
a spanning path anchors a positive fraction of the rows or columns
in a fixed enlargement of a touched cell, giving an \(\ell^2\)
factor; the enlargements have bounded overlap.  Duality gives
\eqref{eq:profile-path-cost}, and \eqref{eq:profile-chain-mgf}
transfers it to the cable posterior.
The local estimate is \cite[Equation~(3.6), Lemma~3.4]{X}.

We also need the finite estimates that control the sigma-finite
constant before band weighting.  Write its sign disintegration as
\(\nu_{n,N}(d\omega)\mathcal R_\omega(d\sigma)\), and put
\[
 M_C(f)=\sum_{x\in C}|h_x|f_x,\qquad
 m_C=M_C(u),\qquad Q=\sum_Cm_C^2,
\]
where \(u\) is the probability weight on all primary vertices of
\(U\).  If \(u_d\) is the probability weight in a cell \(d\),
we abbreviate \(M_C(u_d)\) to \(M_C(d)\).  A cluster is counted once
if it contains a primary vertex;
internal cable vertices contribute no mass.  For \(t\ge1\) and every
integer \(q\ge1\), the bin bounds are
\begin{align}
 \nu_{n,N}(Q\le4t^2)&\le C(1+t),\notag\\
 \int_{Q\le4t^2}\left(\sum_CM_C(f)^2\right)^q\,d\nu_{n,N}
 &\le C_q(1+t)
 \left(t\left|\sum_xf_x\right|
       +\mathcal R\left(f-u\sum_xf_x\right)^{1/2}\right)^{2q}.
 \label{eq:profile-bins}
\end{align}
For any nonnegative cost \(I(C)\) determined by a completed cluster,
and all sufficiently large \(t=t(r)\), reinsertion gives
\begin{equation}\label{eq:profile-reinsertion}
 \int\sum_{\diam C\ge r}I(C)\one_{\{m_C\le2t\}}\,d\nu_{n,N}
 \le2\int\one_{\{Q\le8t^2\}}
          \sum_{\diam C\ge r}I(C)\one_{\{m_C\le2t\}}\,d\nu_{n,N}.
\end{equation}
Finally, for every \(j>0\), some \(M<\infty\) satisfies
\begin{equation}\label{eq:profile-damping}
 \mathcal P_{n,N}\left((1+|S|)^{-M}
             \one_{\{t^2<Q\le4t^2\}}\right)
                       \le C_j(1+t)^{-j}.
\end{equation}
The constants in these estimates depend only on fixed geometry,
the displayed moment orders, \(r\) when present, and \(J\).
Indeed, the proof of \eqref{eq:profile-bins} uses fair-sign
Khintchine and Paley--Zygmund bounds, at most \(C(1+t)\)
translations with \(|S|\le Ct\), and the primary reference moments.
After a macroscopic cluster is exposed, path anchoring bounds every
moment of the complementary squared-mass sum; its probability of
being at most \(4t^2\) is at least \(1/2\), proving
\eqref{eq:profile-reinsertion}.  Polynomial damping follows by
separating clusters below a fixed diameter from the macroscopic
ones: a small cluster misses one of finitely many fixed boxes and
is controlled by the corresponding mean-free test; a large cluster
anchors the complementary mean.  These are finite estimates, proved
in \cite[Lemma~3.5, Equations~(3.7)--(3.9)]{X}; their only parameter-dependent inputs are the
uniform primary moment bounds and \eqref{eq:profile-chain-mgf}.
All cable statements are made first at a fixed primary graph and then
passed through its subdivision limit.  Thus they require no convergence
rate for \(N\) uniform in \(k\).

\emph{Band-weighted clusters and tilt errors.}
We now combine those finite estimates with
\eqref{eq:profile-band-density}.  This is the point where the free
limits enter the exploration argument.  If
\(F=(\sum_CM_C(f)^2)^{1/2}\), or if \(F\) is a signed partial sum
selected using signless data,
and \(A_t\) is a dyadic bin of \(Q\), then
\begin{align*}
 \mathcal P_{n,N}(\widehat W_N|F|^q\one_{A_t})
 &\le
 \mathcal P_{n,N}(\widehat W_N^4(1+|S|)^M)^{1/4}
 \mathcal P_{n,N}((1+|S|)^{-M}\one_{A_t})^{1/4}\\
 &\hspace{35mm}\times
 \mathcal P_{n,N}(|F|^{2q}\one_{A_t})^{1/2}.
\end{align*}
The bin moments grow polynomially in \(t\), whereas
\eqref{eq:profile-damping} has arbitrarily high polynomial decay.
Choosing its exponent after \(q\) therefore bounds the sum and its
tail.  In particular, with \(D_\ell^2=1+\log(1/\ell)\),
\begin{equation}\label{eq:profile-weighted-moments}
 \E_{W,N}\left(\sum_CM_C(d)^2\right)^q
       +\E_{W,N}|d_*|^{2q}
        \le C_{q,\lambda}D_\ell^{2q}+o_k(1).
\end{equation}
Here \(d_*\) is any signed selection whose nonnegative coefficients
are bounded by the full cell average.  The logarithm is the ordinary
reference cost of a probability cell average minus \(S\).  Fixed
observation boxes have a fixed constant in its place.  The remainder
is taken at fixed \(\ell,\lambda\); the constants remain valid in
the subsequent cell refinement.

Replacing the signs of all clusters of diameter at least \(r\) by
independent fair signs changes the band law by vanishing total
variation in the first two limits of \eqref{eq:profile-band-order}.
Here is the quantitative estimate ensuring this assertion for our
sequence.  On \(Q\le t^2\), expand the difference of the two band
energies in the old and new signs.  Fair-sign orthogonality, followed
by exposure of one cluster, \eqref{eq:profile-path-cost}, and
\eqref{eq:profile-reinsertion}, gives
\begin{equation}\label{eq:profile-sign-error}
 \mathcal P_{n,N}\left(\one_{\{Q\le t^2\}}
       |\log W(\sigma')-\log W(\sigma)|^2\right)
       \le C_{\lambda,t,r}\ell^2D_\ell^2+o_k(1).
\end{equation}
The intermediate conditional second moment is bounded by
\(C\ell^2\sum_d|d|M_C(d)^2\), which explains the factor
\(\ell^2\).  On this finite-mass bin, split at a fixed tolerance
for the logarithmic difference, use H\"older and
\eqref{eq:profile-band-density} on the exceptional set, and use
\(|e^x-e^y|\le(e^\eta-1)(e^x+e^y)\) when \(|x-y|\le\eta\).
Then send \(\eta\) to zero.  The weighted moments remove the bin
cutoff, also against fixed polynomial observables.  Exchanging the
old and new sign samples preserves the unweighted measure and \(Q\),
so the two weighted tails obey the same bound.  This supplies the
sequence version of the argument in \cite[Lemmas~3.6--3.7]{X}.

Clusters of diameter less than \(r\) contribute a second moment at
most \(C_t r^2(1+\log(1/r))\) on a fixed bin to a bounded density
smear: partition into \(r\)-cells and use that such a cluster meets
only boundedly many neighboring cells.  Higher bin moments and the
weighted tail bound extend the vanishing to every fixed moment.
It follows, by fair-sign resampling, the band Gaussian limit, and
then removal of smaller clusters, that for nonnegative fixed
density smears \(f_1,f_2\),
\begin{equation}\label{eq:profile-cluster-cross}
 \limsup_{\ell\downarrow0}\limsup_{k\to\infty}\limsup_{N\to\infty}
 \E_{W,N}\sum_{\diam C\ge r}M_C(f_1)M_C(f_2)
                    \le (f_1,T_\lambda^{-1}f_2).
\end{equation}
If \(Q_f=\sum_{\diam C\ge r}M_C(f)^2\), fair-sign Jensen also
bounds the limiting \(q\)-th moment of \(Q_f\) by the Gaussian
\(2q\)-th moment of its smear.  These bounds are uniform for
\(\lambda\ge\lambda_0\), since the Gaussian covariance decreases
with \(\lambda\).  Uniformity toward large strength is thus obtained
after the mesh and cell limits.

During exploration write \(h_d=d_*+h_d^U\).  Relative to its centered
posterior with the restricted band precision, the unvisited law has
linear tilt
\[
 L=-\lambda\sum_d|d|d_*h_d^U,\qquad
 \zeta^2=C\lambda^2\ell^2\sum_d|d|d_*^2.
\]
The latter quantity bounds its reference cost by
\eqref{eq:profile-path-cost}.  For one completed macroscopic cluster
define \(\zeta_C^2\) by replacing \(d_*\) with \(M_C(d)\), and put
\(E_2=\sum_{\diam C\ge r}\zeta_C^2\).  Weighted Jensen and
\eqref{eq:profile-weighted-moments} give
\begin{equation}\label{eq:profile-tilt-moments}
       \E_{W,N}\zeta^{2q}+\E_{W,N}E_2^q
        \le C_{q,\lambda,r}\ell^{2q}D_\ell^{2q}+o_k(1).
\end{equation}
When \(\zeta\le1\), the centered posterior \(\mu\) satisfies
\[
 1\le\mu(e^L),\qquad
 \mu(e^{2L})\le e^{C\zeta^2},\qquad
 \mu((e^L-1)^2)\le C'\zeta^2.
\]
Thus fixed-cost posterior means and second moments change by
\(O(\zeta)\), using their fourth moments.  Large tilts are removed by
polynomial truncation.  In the size-biased exploration used below,
the precise bound for fixed boxes \(B,i\) is
\begin{equation}\label{eq:profile-large-tilts}
 \sum_{\diam C\ge r}M_C(B)\one_{\{\zeta_C>1\}}
       (1+h_i^2+M_C(i)^2)
       \le \sqrt{Q_BE_2}\,(1+h_i^2+Q_i).
\end{equation}
For \(\zeta_C\le1\), the corresponding errors are bounded by
\(\sqrt{Q_B}(1+\sqrt{Q_i})(\sqrt{E_2}+E_2)\).
H\"older, the weighted moments, and \eqref{eq:profile-tilt-moments}
make both expectations vanish at fixed \(\lambda\) after the mesh
and cell limits.  For exploration from a filled region the same
conclusion follows directly by discarding \(\zeta>1\) using
\eqref{eq:profile-tilt-moments} and the actual band moments.
The finite tilt estimates used here are
\cite[Equations~(3.12)--(3.14)]{X}; the argument above identifies
the sequence limits of all their moment inputs.

\emph{Vanishing reach.}
It remains to show that exploration from \(p\) is unlikely to cross
the band and reach \(V\).  Choose finite small-cell coverings
\(i\) and \(B\) on the starting and target sides, respectively,
with fixed positive gaps from the band, the starting region, and
the target region.  They cover separator strips,
so any crossing cluster meets a cell on each side.  The radius
\(\rho\) of these cells is fixed before the band refinement.
Choose the macroscopic threshold \(r>0\) below the fixed separations
that a cluster joining the starting region, either comparison strip,
the band, or the target must span.  Every cluster contributing to
the cross terms below then has diameter at least \(r\), and the
previous moment and tilt estimates apply with this one threshold.
Path anchoring bounds the remaining centered variance of an average
in a reached cell by a constant \(C_0\) independent of \(\rho\).
On the other hand, a logarithmic trial function supported away from
the band gives the continuum band variance lower bound
\[
                   c a_*\log(1/\rho)-C.
\]
The trial pays no band mass, so this bound is uniform in \(\lambda\).
At a free wall it is restricted to the rectangle, which is permitted
for the Neumann form.  Since \(a_*>0\), choose \(\rho\) so small
that the prior variance exceeds \(C_0+2c_0\) for a fixed \(c_0>0\).
The band limits then give a gap at least \(c_0\) in the discrete
laws, for each fixed strength and sufficiently fine partition and mesh.
This is the only step that requires a positive height coefficient;
it uses \(a_*>0\) in the limiting band law.  The finite estimates
require only \(b\in J\), with no assumption on \(a(b_k)\).

There are two useful consequences of this gap.  Write a target
average as exposed plus unvisited height.  The centered posterior
variance is at most the prior variance and loses at least \(c_0\)
if the cell is reached.  Restoring its tilt has the vanishing
errors just proved.  Fair-sign resampling of the exposed term gives
\begin{equation}\label{eq:profile-exterior-deficit}
 c_0\Pbb_{W,N}(B\text{ reached})
    \le\E_{W,N}\sum_{C:C\cap p\ne\varnothing}M_C(B)^2+o(1).
\end{equation}
Here and below \(o(1)\) vanishes in the first two limits at fixed
\(\lambda\).  A second comparison starts by biasing with \(|h_x|\),
for \(x\in B\), and exposing its cluster.  The finite size-bias
inequality
\[
                   \E[X^2|Y|]\ge\E X^2\,\E|Y|
\]
holds for real linear tests under centered band precision.  It follows
from the minimal characteristic logarithmic Hessian in the finite
comparison rules: integrate
\(\E[X^2\cos(tY)]\le\E X^2\,\E\cos(tY)\) against
\(2\,dt/(\pi t^2)\) after subtracting the identity at zero.
Consequently the pre-exposure second moment at \(i\) is at least
its unbiased variance.  The completed-cluster bias cancels from
the conditional complement law.  Apply the same variance deficit,
multiply by \(\E|h_x|\), and average \(x\) over \(B\).  This yields
\begin{equation}\label{eq:profile-biased-deficit}
 c_0\E_{W,N}\sum_{C:C\cap i\ne\varnothing}M_C(B)
       \le\E_{W,N}\sum_C M_C(B)M_C(i)^2+o(1).
\end{equation}
Only clusters spanning the fixed separation contribute to the
exposed or tilted terms.  They have a fixed positive diameter,
so \eqref{eq:profile-large-tilts} justifies their summed errors.
Equations \eqref{eq:profile-exterior-deficit} and
\eqref{eq:profile-biased-deficit} are the two finite deficit
comparisons in \cite[Equations~(3.16)--(3.17),
proof of Proposition~3.8]{X}.

As \(\lambda\to\infty\), the inverse forms \(T_\lambda^{-1}\)
decrease to the inverse with a hard zero band on \(K\).  This follows
from weak energy compactness and recovery by functions zero on \(K\).
The band separates \(B\) from \(i\), so
\((B,T_\lambda^{-1}i)\to0\).  By
\eqref{eq:profile-cluster-cross}, the limiting expectation of
\(\sum_CM_C(B)M_C(i)\) therefore vanishes.  On
\(\max_CM_C(i)\le R\), the right side of
\eqref{eq:profile-biased-deficit} is at most \(R\) times this cross
sum.  Its remaining tail tends to zero as \(R\to\infty\), uniformly
after the first two limits and toward large \(\lambda\), because
\[
              \sum_CM_C(B)M_C(i)^2\le\sqrt{Q_B}\,Q_i
\]
and the \(Q\)'s have arbitrary fixed moments uniformly in that
order of limits.  Thus the first-mass sum on the left of
\eqref{eq:profile-biased-deficit} vanishes.  Every cluster joining
\(p\) to \(B\) visits some cell \(i\).  Truncating
\(\max_CM_C(B)\) in \eqref{eq:profile-exterior-deficit} and applying
these first-mass bounds, with the same moment control of its tail,
proves
\begin{equation}\label{eq:profile-vanishing-reach}
 \lim_{\lambda\to\infty}\limsup_{\ell\downarrow0}
 \limsup_{k\to\infty}\limsup_{N\to\infty}
 \Pbb_{W,N}\bigl(\text{exploration from }p
                     \text{ reaches }V\bigr)=0.
\end{equation}

\emph{Pinned tests and the grounded approximation.}
On a finite subdivided graph, explore every cluster meeting the
primary pin region.  If it does
not reach \(V\), the exposed part of \(X_{g,n}\) is zero.
The centered posterior has the pin \(p\), possibly further zero
pins and increased segment precision, and a nonnegative restricted
band form.  Its variance is at most the variance of \(X_{g,n}\)
under the primary hard-pin law on that subdivided graph, by finite
precision comparison.  At fixed primary graph, this latter variance
converges as \(N\to\infty\) to the Bessel hard-pin variance.
The averaged tilt errors vanish as above.  On failure,
Cauchy--Schwarz bounds the actual second-moment contribution by
the fourth-moment root times the square root of the failure
probability.  The band fourth moment is uniformly bounded in the
successive limits, whereas \eqref{eq:profile-vanishing-reach}
makes that probability vanish.  Hence
\[
 \liminf_{k\to\infty}\Var_p X_{g,n}
                   \ge a_*\mathcal E_U(g,g).
\]
The free variance bound gives the opposite limsup.  For every real
\(t\), the finite lower bound for the logarithmic moment generating
function and the upper bound obtained by releasing the pin give
\[
 \frac{t^2}{2}\Var_p X_{g,n}
       \le \log\E_p e^{tX_{g,n}}
       \le \log\E_{\rm free}e^{tX_{g,n}}.
\]
The free release is permitted because \(A_{U,n}g_n\) has total zero.
The two endpoints converge to \(a_*t^2\mathcal E_U(g,g)/2\).
Applying this to finite linear combinations gives the joint
Gaussian Laplace limit for these tests; uniform primary exponential
moments give their polynomial moment limits.  Weak free-face terms
may be approximated first by bounded density smears in primary
reference cost, with the gap from \(p\) and \(K\) preserved.

Finally let \(f\) be any bounded piecewise continuous density,
and set \(v=G_pf\).  Smooth shifts \(g_j\) vanishing near \(p\),
and allowed to meet the free walls, are dense in \(\mathcal H_p\)
in energy.  Choose \(\mathcal E_U(v-g_j,v-g_j)\to0\).
The ordinary mixed-form approximation on these rectangular
geometries, as stated in \cite[Section~3.1 and proof of
Theorem~3.1]{X}, gives
\begin{equation}\label{eq:profile-grounded-approximation}
 \lim_{k\to\infty}
 \mathcal R_{p,n}\bigl(f_n-A_{U,n}g_{j,n}\bigr)
                    =\mathcal E_U(v-g_j,v-g_j),
\end{equation}
where \(\mathcal R_{p,n}\) is the squared dual norm of the unit-edge
form grounded on \(p\).  To see the precise form of this approximation,
let \(A_{p,n}\) be the restriction to unpinned coordinates and put
\(v_n=A_{p,n}^{-1}f_n\).  The left side before its limit is
\[
 (f_n,A_{p,n}^{-1}f_n)-2(f_n,g_{j,n})
                +(g_{j,n},A_{U,n}g_{j,n}).
\]
Mixed inverse-form convergence, quadrature, and smooth energy
consistency give respectively
\(\mathcal E_U(v,v)\), \(\mathcal E_U(v,g_j)\), and
\(\mathcal E_U(g_j,g_j)\), proving
\eqref{eq:profile-grounded-approximation}.  Values of a test on pinned vertices are
immaterial in this norm.  Equivalently, a flow representing its
unpinned coefficients may discharge into the pinned vertices.
This is why \eqref{eq:profile-grounded-approximation} imposes no
zero-total condition on \(f\).

Uniform reference-cost bounds make the difference between the two
smears in \eqref{eq:profile-grounded-approximation} tend to zero
in every fixed moment, first as \(k\to\infty\) and then
\(j\to\infty\).  Their fixed linear exponential bounds give the
same passage for Laplace transforms, by Cauchy--Schwarz applied to
\(e^{t(h,f_n)}-e^{tX_{g_j,n}}\).  The limiting variance is
\(a_*\mathcal E_U(v,v)=a_*\int_UfG_pf\).
Apply the argument to every finite linear combination of the
\(f_i\)'s and use polarization to obtain
\eqref{eq:profile-pinned-covariance}, with the asserted joint
Laplace and moment convergence.
\end{proof}

Two choices of pin will be used below.  If \(D\) is an inner rectangle
and \(p=U\setminus D\) is its positive-width rectangular rim, then
\(\mathcal H_p\) identifies
with \(H^1_0(D)\) by zero extension, and
\eqref{eq:profile-pinned-covariance} is the Dirichlet covariance
\(a_*G_D\).  If a positive-width primary collar surrounds a compactly
supported shift inside an unwrapped torus rectangle, zero pinning
that collar separates its interior law from the exterior.  The lemma
therefore applies to that interior law as well.


\subsection{Torus limits and uniform zero-coefficient estimates}

\begin{proposition}[Gaussian laws at continuity scales]\label{prop:gaussian-profile}
Let $w$ be a continuity point of the extracted profile, let
$\delta_k\log_Ln_k\to w$, and put $a_*=\mathcal A(w)$.
The free rectangular limits are those in
Lemma~\ref{lem:profile-free}; if $a_*>0$, their fractional means have
the conclusion of Lemma~\ref{lem:profile-phase}, and their fixed
primary pinned geometries have the Gaussian Laplace and moment limits
of Lemma~\ref{lem:profile-pins}. In particular, on a square interior
surrounded by a fixed positive-width zero rim, the covariance of
bounded piecewise continuous density smears converges to
$a_*(-\Delta_D)^{-1}$, with
Dirichlet boundary on that square.

On the zero-winding height torus of side $n_k$, every smooth periodic
shift $g$ satisfies
\begin{equation}\label{eq:profile-torus}
 \Var(h,A_{\mathrm{tor}}g_k)\longrightarrow
                         a_*\int_{\mathbb T^2}|\nabla g|^2.
\end{equation}
Thus the area-weighted first cosine observation in
\eqref{eq:physical-cosine} has limiting variance
$a_*\mathfrak g$, where $\mathfrak g=1/(8\pi^2)$.
\end{proposition}

\begin{proof}
The free and pinned assertions have been proved. For the torus upper
bound, cut into small fixed rectangles, restrict the gradient test to
each piece, and use \eqref{eq:profile-partial-upper}. The omitted seams
and the gradient approximations have uniformly vanishing reference
cost. Therefore the limsup in \eqref{eq:profile-torus} is at most its
right side. This already proves the assertion when $a_*=0$.

For $a_*>0$, first suppose that $g$ is supported in an unwrapped
rectangle strictly inside one torus chart. Add primary zero pins on a
collar of fixed positive width surrounding its support inside that
chart. Centered precision comparison decreases variance. The collar
separates the interior law from the torus exterior, and
Lemma~\ref{lem:profile-pins} gives the interior variance limit.
Indeed the torus test equals the grounded interior test there, and
$(-\Delta_D)^{-1}(-\Delta g)=g$ because $g$ vanishes near the
collar. Its variance is therefore $a_*\int|\nabla g|^2$, matching
the upper bound. For general $g$,
use a finite smooth partition of unity to write $g$ as a sum of shifts
supported in such charts. On each joint covariance subsequence the
upper-bound defect is a positive semidefinite quadratic form. Since
its diagonal vanishes on each summand, its cross terms with those
summands vanish as well. It therefore vanishes on their sum.

For the cosine density $f(x)=\cos(2\pi x_1)$, its inverse torus
Laplacian is $g=f/(4\pi^2)$ and
$\int|\nabla g|^2=1/(8\pi^2)$. The discrete eigenvalue converges to
$4\pi^2$ after multiplication by $n_k^2$. The normalizing correction
has vanishing reference cost, which proves the last assertion.
\end{proof}

\begin{corollary}[Uniform affine transforms on a zero interval]
\label{cor:profile-zero-uniform}
Suppose $I=[u,v]\subset(0,\infty)$ consists of continuity points of
$\mathcal A$ and $\mathcal A=0$ on $I$. Fix a compact set
$K\subset\R^2$ and $\lambda\in\R$. For a free square of side $m$
write $X_{m,p}=(h,A_{B_m}(p\cdot x/m))$. Then
\begin{equation}\label{eq:profile-zero-uniform}
 \sup_{\substack{m\in\N:\ \delta_k\log_Lm\in I\\p\in K}}
 \left|\E_{B_m,b_k}^{\mathrm{free}}e^{\lambda X_{m,p}}-1\right|
                                      \longrightarrow0.
\end{equation}
The assertion is unchanged by translating the squares or multiplying
their sides by factors in a fixed compact subset of $(0,\infty)$,
provided their resulting logarithmic scales stay in a fixed
zero-profile neighborhood of $I$.
\end{corollary}

\begin{proof}
If $a_m(b_k)$ failed to tend to zero uniformly on the stated set,
there would be offending $m_k$ and a subsequence with
$\delta_k\log_Lm_k\to w\in I$. This contradicts
\eqref{eq:profile-limit}. Square symmetry gives
$\Var X_{m,p}=a_m(b_k)|p|^2$, uniformly small for $p\in K$.
The reference exponential bound at $\pm2\lambda$ supplies uniform
integrability of $e^{\lambda X_{m,p}}$. Since $X_{m,p}\to0$ in
probability uniformly over these choices, the expectations tend to
one uniformly. Translation does not change the law. The extra side
factors change $\delta_k\log_Lm$ by $o(1)$; compactness gives the
last statement.
\end{proof}

\subsection{Excluding intermediate positive coefficients}

The fixed-temperature coefficient gap in \cite[Theorem~2.5]{X}
is proved by growth
of a face characteristic coefficient. For a moving temperature, a
single growth comparison cannot be iterated at arbitrarily many
scales. We instead integrate those comparisons over a compact
logarithmic interval. The initial characteristic coefficient must
therefore have a quantitative lower bound on this interval.

Let $\rho_m^{\mathrm{sym}}$ be one quarter of the sum of the four
face probabilities of $B_m$, and define
\begin{equation}\label{eq:profile-z}
 z_m(b)=\E_{B_m,b}^{\mathrm{free}}
                           e^{i(h,\rho_m^{\mathrm{sym}})}.
\end{equation}
It is strictly positive and at most one.

\begin{lemma}[Relative growth of the face phase]\label{lem:profile-ratio}
There is $c>0$, independent of sufficiently large fixed integer $G$,
with the following property. If $w$ is a continuity point and
$\mathcal A(w)=a_*>0$, then along every sequence
$\delta_k\log_Lm_k\to w$,
\begin{equation}\label{eq:profile-ratio}
 \liminf_k\frac{z_{Gm_k}(b_k)}{z_{m_k}(b_k)}
                          \ge cG^{2-a_* /(4\pi)}.
\end{equation}
For each fixed such $G$ there is also $c(G)>0$ such that
\begin{equation}\label{eq:profile-ratio-rough}
                     z_{Gm}(b)/z_m(b)\ge c(G)
\end{equation}
for all sufficiently large $m$ and all $b$ in the fixed compact
temperature interval.
\end{lemma}

\begin{proof}
We use the finite soft joining formula
\cite[Equation~(2.11)]{X}. Tile $B_{Gm}$ by $G^2$ free
$m$-squares. The outer-face phase has cell totals $\xi_i$, supported
on boundary cells, with $\sum_i\xi_i=1$ and
$\sum_i\xi_i^2\le C/G$. On the space of neutral cell totals choose a
flow right inverse $P$ on the neighboring-cell graph, and let
$\rho_{ij}$ be the probability on the face of cell $i$ adjoining
cell $j$. Define a vector
$\mathsf H$ in that neutral Euclidean space by
\[
 (\mathsf H,d)=\sum_{\langle i,j\rangle}P(d)_{ij}
       \bigl((h_i,\rho_{ij})-(h_j,\rho_{ji})
                         -\bar\eta_{ij}\bigr),
\]
where $\bar\eta_{ij}$ is the mean of the independent Bessel seam
increments with the same orientation. Translating cells by constants
changes $\mathsf H$ by the neutral projection of those constants. Penalize by
$e^{-|\mathsf H|^2/2}$ and sum relative translations, quotienting by their
common translation. The auxiliary-confinement argument in the proof
of Lemma~\ref{lem:profile-phase} justifies the finite precision
comparisons for these sums.
Poisson summation gives the positive numerator
\begin{equation}\label{eq:profile-soft-sum}
 \sum_{\substack{q_i\in\Z\\\sum_iq_i=1}}
 e^{-|q-\xi|^2/2}\,
 \prod_i\phi_{m,b}\left(s_i+
              \sum_{j\sim i}P(q-\xi)_{ij}\rho_{ij}\right)
 \chi_m(P(q-\xi)),
\end{equation}
up to a common scalar. Here $s_i$ is the restriction of the outer
phase to cell $i$, $\rho_{ij}$ is its seam-face probability, and
$\chi_m$ is the characteristic function of the independent seam
averages. The denominator is the corresponding sum with $s=0$ and
total zero. This follows by Fourier transforming the Gaussian penalty
on the neutral space. The reciprocal of the projected cell-translation
lattice is the neutral integer lattice. Exact seam matching adds centered precision, so the joined
coefficient is at least this soft coefficient.

For fixed $G$, the denominator tends to one along the sequence in
\eqref{eq:profile-ratio}. Its zero sector is one. Every nonzero
sector contains a nonzero integer cell total and tends to zero by
Lemma~\ref{lem:profile-phase}; the Gaussian weights permit dominated
convergence. Each fixed-flux seam characteristic tends to one because
the seam averages have variance $O(m^{-1})$.

Retain the sectors $q=e_{i_0}$ with $i_0$ at distance at least $G/4$
from the boundary. The radial flow of \cite[Equation~(2.14)]{X} has outward face
fluxes $H_{i,\pm j}$, constant outer flux $1/(4G)$, and cell
divergence $e_{i_0}$. Set
\[
 D_{ij}=\tfrac12(H_{i,+j}-H_{i,-j}),\qquad
 S_{ij}=H_{i,+j}+H_{i,-j}.
\]
The unit-cell radial field, corrected to have constant boundary flux,
gives
\begin{equation}\label{eq:profile-radial-cost}
 \sum_{i,j}D_{ij}^2\le\frac1{2\pi}\log G+C,
 \qquad \sum_{i,j}S_{ij}^2\le C.
\end{equation}
One may obtain it by integrating
$(x-x_0)/(2\pi|x-x_0|^2)+E(x)$ across cell faces, where the Neumann
boundary correction has $|E|\le C/G$ and $|\nabla E|\le C/G^2$.
The radial energy is $(2\pi)^{-1}\log G+O(1)$, and the face-sum
errors are summable. The constants are uniform for the retained
central cells. Prescribe $P(e_{i_0}-\xi)$ to be these flows for all
retained $i_0$ and extend linearly. This is possible because the
vectors $e_{i_0}-\xi$ are linearly independent, their central
coordinates being distinct.

In each cell expand the logarithm of its characteristic coefficient
about the symmetric four-face profile of the same total. Square
symmetry makes the linear term vanish on every zero-total face
perturbation. The logarithmic Hessian bound thus gives the product
lower bound
\begin{equation}\label{eq:profile-relative-lower}
        z_m(b)\exp\left\{-\tfrac12\sum_i\Var(h,Y_i)\right\},
\end{equation}
where $Y_i$ is the difference from that symmetric profile. Since
$q_i=S_{i1}+S_{i2}$, its decomposition is explicitly
\begin{align*}
 Y_i={}&\sum_{j=1}^2D_{ij}(\rho_{i,+j}-\rho_{i,-j})\\
 &+\frac{S_{i1}-S_{i2}}4
   (\rho_{i,+1}+\rho_{i,-1}-\rho_{i,+2}-\rho_{i,-2}).
\end{align*}
Square reflections make these three components mutually orthogonal.
The two odd face differences have variance exactly $a_m(b)$; the even
component has reference cost at most $C(S_{i1}^2+S_{i2}^2)$. Lemma~\ref{lem:profile-free} and
the reference bound give
\[
 \limsup_k\sum_i\Var(h,Y_i)
       \le a_*\sum_{i,j}D_{ij}^2+C\sum_{i,j}S_{ij}^2.
\]
The Gaussian sector weight is bounded below because
$|e_{i_0}-\xi|$ is bounded. There are at least $c_0G^2$ such sectors.
Combining these facts with \eqref{eq:profile-radial-cost} proves
\eqref{eq:profile-ratio}. The additive energy constant and the even
component constant are uniform, since $a_*\le a$ and the reference
bounds are uniform. Therefore its prefactor $c$ does not depend on
$G$ or on the positive continuity value.

For \eqref{eq:profile-ratio-rough}, keep $G$ fixed and retain just one
central sector. Bound every variance in
\eqref{eq:profile-relative-lower} by its reference cost. The soft
denominator is at most the fixed finite-dimensional Gaussian lattice
sum, since all its characteristic factors are at most one. The
selected seam factor has a uniform positive lower bound: a mean over
$m$ independent Bessel increments at fixed flux $q$ has
characteristic function
$\exp\{bm(\cos(2\pi q/m)-1)\}$, which is bounded below uniformly
in large $m$ and in the compact temperature interval. There are only
finitely many seams. These bounds prove \eqref{eq:profile-ratio-rough}.
Crucially, \eqref{eq:profile-relative-lower} retains the factor
$z_m(b)$ regardless of how small it is.
\end{proof}

\begin{lemma}[Critical lower bound for the face phase]
\label{lem:profile-subexponential}
For the fixed buffer $D_0$ and block ratio $L$,
\begin{equation}\label{eq:profile-subexponential}
 -\log z_{D_0L^j}(\bc)=o(j),\qquad j\to\infty.
\end{equation}
Moreover $z_m(b)\ge z_m(\bc)$ whenever $b\le\bc$.
\end{lemma}

\begin{proof}
At $\bc$, \cite[Equation~(2.15)]{X}, the fixed-temperature version
of \eqref{eq:profile-ratio}, has coefficient $a=8\pi$. Taking $c_1$ to be the smaller of $1$ and half the prefactor in that
liminf bound, we obtain a constant independent of sufficiently large
fixed $r$ such that
for $G=L^r$,
\[
 z_{Gm}(\bc)\ge c_1z_m(\bc)
\]
for all $m$ beyond a threshold depending on $r$. Fix $r$ and iterate
on each residue class of $j$ modulo $r$. The finitely many initial
coefficients are strictly positive. Hence
\[
 \liminf_{j\to\infty}j^{-1}\log z_{D_0L^j}(\bc)
                                      \ge r^{-1}\log c_1.
\]
Now let $r\to\infty$ and use $z_m\le1$ to obtain
\eqref{eq:profile-subexponential}. Characteristic coefficients
decrease with $b$ by the chain precision comparison, giving the last
claim. The lower bound is therefore uniform for all $b\le\bc$.
\end{proof}

\begin{proposition}[The profile has one jump, at the passage time]
\label{prop:profile-step}
The profile extracted in Proposition~\ref{prop:profile} satisfies
\begin{equation}\label{eq:profile-step}
 \mathcal A(w)=a\quad(0<w<T),\qquad
 \mathcal A(w)=0\quad(w>T),
 \qquad T=\frac{\pi}{\sqrt{3\mathcal P}}.
\end{equation}
In particular all points of these two open intervals are continuity
points. No assertion about its value at the single point $T$ is
needed.
\end{proposition}

\begin{proof}
First we exclude a continuity value in $(0,a)$. If such a value
exists, monotonicity and continuity give a compact interval
$[u,v]\subset(0,\infty)$ of positive length and $\eta>0$ such that
all its continuity values lie in $[\eta,a-\eta]$. Choose a fixed
integer stride $r$ so large that, for $G=L^r$,
$cG^{2-(a-\eta)/(4\pi)}>2$, where $c$ is from
Lemma~\ref{lem:profile-ratio}.

On one residue class modulo $r$ take consecutive indices $j$ for
which $\delta_k j$ runs across $[u,v]$, and define on their cells
of width $\delta_kr$ the step function
\[
 F_k(t)=\log\frac{z_{D_0L^{j+r}}(b_k)}{z_{D_0L^j}(b_k)}.
\]
Choose the endpoints of this union of cells to converge to $u,v$.
At almost every $t\in(u,v)$ the profile is continuous. The starting
cell size has logarithmic scale tending to $t$, so
\eqref{eq:profile-ratio} gives $\liminf_k F_k(t)\ge\log2$.
Equation~\eqref{eq:profile-ratio-rough} gives the common lower bound
$F_k\ge\log c(G)$. Extend by zero outside the chosen cells in a
fixed enclosing interval. Fatou's lemma therefore gives
\begin{equation}\label{eq:profile-fatou}
          \liminf_k\int F_k(t)\,dt\ge(v-u)\log2>0.
\end{equation}
On the other hand, the integral telescopes exactly:
\[
 \int F_k(t)\,dt
 =\delta_kr\bigl[\log z_{D_0L^{j_+(k)}}(b_k)
                    -\log z_{D_0L^{j_-(k)}}(b_k)\bigr].
\]
Since $z\le1$, Lemma~\ref{lem:profile-subexponential} bounds its
limsup by
$\limsup_k\delta_kr\,o(j_-(k))=0$; here
$\delta_kj_-(k)\to u>0$. This contradicts
\eqref{eq:profile-fatou}. Thus every continuity value belongs to
$\{0,a\}$.

For $0<w<T$, Proposition~\ref{prop:passage} gives the torus cosine
variance limit $a\mathfrak g$, whereas
Proposition~\ref{prop:gaussian-profile} gives
$\mathcal A(w)\mathfrak g$ at continuity points. Thus
$\mathcal A(w)=a$ at all such points. Suppose a continuity point
$w_0>T$ also had value $a$. Monotonicity and the bound
$\mathcal A\le a$ would make the profile identically $a$ on a
neighborhood of $T$, so $T$ would be a continuity point. The moving
exit sizes $n_{j_e}=D_0L^{j_e}$ in Proposition~\ref{prop:passage}
have $\delta_kj_e\to T$. Proposition~\ref{prop:gaussian-profile}
would give cosine variance tending to $a\mathfrak g$ there,
contradicting the strictly smaller exit limsup. Every continuity
point above $T$ therefore has value zero. Finally a monotone function
whose values on the continuity points of an open interval are a
single constant is that constant everywhere on the interval.
This proves \eqref{eq:profile-step}.
\end{proof}
